\section{Therapeutic Strategies Targeting CAFs}
\label{sec:therapy}

The recognition of CAFs as active contributors to tumour progression has spurred the development of therapeutic strategies aimed at eliminating, reprogramming, or exploiting these cells. This section reviews the current landscape of CAF-targeted therapies in breast cancer, organized by mechanism of action.

\subsection{CAF depletion strategies}

The most direct approach to targeting CAFs is their selective elimination, exploiting surface markers or metabolic dependencies that distinguish them from normal fibroblasts.

\subsubsection{FAP-targeted therapies}

Fibroblast activation protein (FAP) is a serine protease highly expressed on CAFs but largely absent from normal fibroblasts \citep{garin2010}. This differential expression has made FAP an attractive therapeutic target.

\textbf{FAP-targeted antibodies:} Sibrotuzumab, a humanized anti-FAP antibody, showed modest activity in early-phase clinical trials but failed to demonstrate significant efficacy as monotherapy \citep{scott2003}. More recently, bispecific antibodies targeting both FAP and immune checkpoints (e.g., FAP $\times$ PD-1) have shown enhanced activity in preclinical models by simultaneously depleting CAFs and reactivating T cells \citep{fang2023}.

\textbf{FAP-targeted CAR-T cells:} CAR-T cells engineered to recognize FAP have shown efficacy in preclinical models of breast cancer, reducing CAF density, disrupting the stromal barrier, and enhancing immune cell infiltration \citep{wang2014}. However, clinical translation has been limited by on-target, off-tumor toxicity, as FAP is expressed at low levels on normal tissues including bone marrow stromal cells and pancreatic stellate cells.

\textbf{FAP-targeted radioligand therapy:} Radiolabeled FAP inhibitors (FAPIs)---originally developed for diagnostic imaging---have shown therapeutic potential when conjugated to beta-emitting radionuclides ($^{90}$Y, $^{177}$Lu) \citep{kratochwil2019}. In breast cancer, FAPI-based radiotherapy has shown promising responses in early-phase trials, with CAF depletion accompanied by tumour regression and immune activation.

\subsubsection{Anti-CAF antibody--drug conjugates}

Antibody--drug conjugates (ADCs) targeting CAF-specific antigens represent another depletion strategy. ADCs against podoplanin (PDPN), a glycoprotein expressed on a subset of CAFs, have shown efficacy in preclinical models \citep{ueji2023}. These agents deliver cytotoxic payloads specifically to CAFs, sparing normal fibroblasts and other cell types.

However, the heterogeneity of CAF marker expression poses a challenge for depletion strategies: no single marker is expressed by all CAF subpopulations, and targeting one marker may spare CAFs expressing alternative phenotypes.

\subsection{CAF reprogramming strategies}

Rather than eliminating CAFs, reprogramming approaches aim to convert tumour-promoting CAFs into quiescent or anti-tumorigenic states.

\subsubsection{All-trans retinoic acid}

All-trans retinoic acid (ATRA) has been shown to reverse CAF activation in multiple cancer types, including breast cancer \citep{liu2023}. ATRA inhibits TGF-$\beta$-induced SMAD signalling, downregulates $\alpha$-SMA expression, and reduces ECM production, effectively converting activated myCAFs into quiescent fibroblasts.

In preclinical models, ATRA combined with chemotherapy enhanced drug delivery by reducing stromal density and improving vascular perfusion \citep{liu2023}. Clinical trials of ATRA in combination with nab-paclitaxel are ongoing in metastatic breast cancer, with preliminary evidence of enhanced anti-tumour activity.

\subsubsection{Vitamin D receptor agonists}

Vitamin D receptor (VDR) signalling opposes fibroblast activation through inhibition of Hedgehog and TGF-$\beta$ pathways \citep{shen2023}. VDR agonists (calcipotriol, paricalcitol) have shown efficacy in reprogramming CAFs in pancreatic cancer, and early evidence suggests similar activity in breast cancer.

In breast cancer models, VDR agonists reduce CAF-derived ECM production, decrease tissue stiffness, and enhance immune cell infiltration \citep{shen2023}. These effects are accompanied by improved response to chemotherapy and immunotherapy, supporting the concept that CAF reprogramming can sensitize tumours to conventional treatments.

\subsubsection{BET inhibitors}

Bromodomain and extra-terminal domain (BET) proteins (BRD2/3/4) are epigenetic regulators that maintain the activated CAF phenotype. BET inhibitors (JQ1, I-BET762) disrupt CAF activation by suppressing the transcription of fibrosis-associated genes \citep{erve2024}.

In breast cancer, BET inhibitors have shown dual effects: they directly suppress tumour cell proliferation while simultaneously reprogramming CAFs toward a quiescent state. This dual activity makes BET inhibitors attractive candidates for combination therapy.

\subsection{Targeting CAF-derived signalling}

Blocking specific signalling pathways downstream of CAFs represents a more targeted approach than global CAF depletion or reprogramming.

\subsubsection{CXCR4/CXCL12 axis}

The CXCR4/CXCL12 axis mediates multiple CAF functions including immune cell recruitment, angiogenesis, and metastasis. CXCR4 antagonists (plerixafor/AMD3100) block these functions and have shown efficacy in preclinical models of breast cancer \citep{domanska2012}.

In clinical trials, plerixafor has shown modest single-agent activity but enhanced efficacy when combined with chemotherapy or immunotherapy \citep{domanska2012}. The combination of plerixafor with anti-PD-1 therapy is being evaluated in clinical trials for metastatic TNBC.

\subsubsection{IL-6/STAT3 axis}

IL-6, primarily secreted by iCAFs, activates STAT3 signalling in both tumour cells and immune cells, promoting proliferation, survival, and immune evasion. Anti-IL-6 receptor antibodies (tocilizumab) and JAK/STAT3 inhibitors are being evaluated in breast cancer \citep{chang2014}.

In preclinical models, IL-6 blockade reduces iCAF-mediated immunosuppression and enhances the efficacy of anti-PD-1 therapy \citep{erve2024}. Clinical trials of tocilizumab in combination with pembrolizumab are underway in metastatic TNBC.

\subsubsection{TGF-$\beta$ signalling}

TGF-$\beta$ is the master regulator of CAF activation and a key mediator of immune suppression. Multiple TGF-$\beta$-targeting strategies are in clinical development:

\textbf{Anti-TGF-$\beta$ antibodies:} Bintrafusp alfa, a bifunctional fusion protein targeting both PD-L1 and TGF-$\beta$, showed promising activity in early-phase trials but failed to meet its primary endpoint in a phase III trial in cervical cancer \citep{strauss2023}. Ongoing trials are evaluating this agent in breast cancer subtypes with high TGF-$\beta$ signalling activity.

\textbf{TGF-$\beta$ receptor kinase inhibitors:} Small-molecule inhibitors of TGF-$\beta$ receptor I (galunisertib, vactosertib) block CAF activation at the receptor level \citep{herbertz2015}. These agents have shown activity in preclinical models and are being evaluated in early-phase clinical trials.

\textbf{Anti-TGF-$\beta$ traps:} ATE-Trap, a TGF-$\beta$ decoy receptor, sequesters TGF-$\beta$ in the extracellular space, preventing it from activating fibroblasts and immune cells \citep{teicher2021}. This approach has shown synergy with anti-PD-1 therapy in preclinical models of breast cancer.

\subsection{Targeting CAF metabolism}

CAFs exhibit distinct metabolic features that can be therapeutically exploited:

\textbf{Glycolysis inhibition:} CAFs rely heavily on aerobic glycolysis (Warburg effect), producing lactate that is exported and consumed by tumour cells \citep{martinez2020}. Inhibiting glycolysis in CAFs (e.g., with 2-deoxyglucose) reduces lactate production and impairs tumour cell metabolism.

\textbf{Glutamine metabolism:} CAFs are major consumers of glutamine, which they convert to glutamate and other metabolites that support tumour growth \citep{bertero2019}. Glutaminase inhibitors (CB-839) can block this metabolic support and have shown activity in preclinical breast cancer models.

\textbf{Lipid metabolism:} CAFs exhibit altered lipid metabolism, with increased fatty acid synthesis and oxidation \citep{sahai2020framework}. Inhibitors of fatty acid synthase (FASN) or carnitine palmitoyltransferase 1 (CPT1) can disrupt CAF metabolic function and impair their tumour-supportive capabilities.

\subsection{Combination strategies}

Given the complexity of CAF biology, combination approaches targeting multiple CAF functions simultaneously may be more effective than single-agent strategies:

\begin{itemize}
  \item \textbf{CAF depletion + immunotherapy}: Combining FAP-targeted CAR-T cells with anti-PD-1 therapy has shown enhanced efficacy in preclinical models by simultaneously removing the physical barrier and reactivating T cells.
  \item \textbf{CAF reprogramming + chemotherapy}: ATRA combined with nab-paclitaxel enhances drug delivery by reducing stromal density while simultaneously reducing CAF-mediated chemoresistance.
  \item \textbf{Anti-TGF-$\beta$ + anti-PD-1}: Blocking TGF-$\beta$ signalling reduces both CAF activation and immune suppression, potentially converting immune-excluded tumours into immune-inflamed tumours.
  \item \textbf{Metabolic inhibitors + immunotherapy}: Targeting CAF metabolism can reduce immunosuppressive metabolite production, enhancing the efficacy of immune checkpoint blockade.
\end{itemize}

These combination strategies reflect the understanding that CAFs contribute to tumour progression through multiple, redundant mechanisms, and that effective therapy must address this complexity.
